**Theorem.** The set of natural numbers equal to their own digit factorial sum is exactly $\{1, 2, 145, 40585\}$. That is,
$$\{\, n \in \mathbb{N} \mid n = \sigma(n) \,\} = \{1,\, 2,\, 145,\, 40585\},$$
where $\sigma(n)$ denotes the sum of the factorials of the decimal digits of $n$.

*Proof.* We prove the two inclusions separately.

$(\subseteq)$ Let $n \in \mathbb{N}$ satisfy $n = \sigma(n)$. By \textbf{Upper Bound}, every fixed point of $\sigma$ satisfies $n \leq 2540160$. Since $n$ is bounded, \textbf{Finite Verification} applies and yields that $n = \sigma(n)$ holds if and only if $n \in \{1, 2, 145, 40585\}$; in particular, $n$ belongs to the right-hand set.

$(\supseteq)$ Let $n \in \{1, 2, 145, 40585\}$. Each such element trivially satisfies $n \leq 2540160$, so \textbf{Finite Verification} again applies. It confirms that every element of $\{1, 2, 145, 40585\}$ is indeed a fixed point of $\sigma$, i.e.\ $n = \sigma(n)$. $\blacksquare$